% ============================================================
%  MTH636 – Game Theory | Lecture 15
%  Proof of Nash's Theorem (Contd.) & Extensive Form Games Intro
% ============================================================
\documentclass[12pt, a4paper]{article}

% ============================================================
%  MTH636 -- Game Theory
%  Notes Preamble — Color-Coded, TikZ-Enhanced
% ============================================================

% --- Core packages ---
\usepackage[T1]{fontenc}
\usepackage[utf8]{inputenc}
\usepackage{lmodern}
\usepackage[margin=2.5cm, top=2.8cm, bottom=2.8cm]{geometry}
\usepackage{microtype}

% --- Math ---
\usepackage{amsmath, amssymb, amsthm, mathtools}
\usepackage{bm}        % bold math
\usepackage{cancel}    % strike-through in math

% --- Color & Boxes ---
\usepackage[dvipsnames, svgnames, x11names, table]{xcolor}
\usepackage[most]{tcolorbox}
\tcbuselibrary{breakable, skins, theorems}

% --- TikZ & Diagrams ---
\usepackage{tikz}
\usepackage{tikz-qtree}
\usetikzlibrary{arrows.meta, positioning, shapes.geometric, fit, backgrounds,
                decorations.pathreplacing, calc, matrix, trees}
\usepackage{pgfplots}
\pgfplotsset{compat=1.18}

% --- Tables ---
\usepackage{booktabs}
\usepackage{multirow}
\usepackage{array}
\usepackage{tabularx}
\usepackage{makecell}

% --- Typography & Layout ---
\usepackage{enumitem}
\usepackage{parskip}
\usepackage{fancyhdr}
\usepackage{titlesec}
\usepackage{titletoc}
\usepackage{hyperref}
\usepackage{graphicx}
\usepackage{float}
\usepackage{caption}

% --- Color Palette ---
\definecolor{defblue}{HTML}{1A5276}      % Definition header
\definecolor{defbluebg}{HTML}{D6EAF8}    % Definition background
\definecolor{thmgreen}{HTML}{1E8449}     % Theorem header
\definecolor{thmbg}{HTML}{D5F5E3}        % Theorem background
\definecolor{exorange}{HTML}{B7770D}     % Example header
\definecolor{exbg}{HTML}{FEF9E7}         % Example background
\definecolor{warnred}{HTML}{922B21}      % Warning/Key Insight header
\definecolor{warnbg}{HTML}{FADBD8}       % Warning background
\definecolor{noticegray}{HTML}{2C3E50}   % Note header
\definecolor{noticebg}{HTML}{EBF5FB}     % Note background
\definecolor{taketawayyellow}{HTML}{7D6608} % Key Takeaway header
\definecolor{takeawaybg}{HTML}{FDFEFE}   % Takeaway background
\definecolor{accentpurple}{HTML}{6C3483} % Accent / proof
\definecolor{lightgray}{HTML}{F2F3F4}    % shading

% ============================================================
%  TCOLORBOX ENVIRONMENTS
% ============================================================

% -- Definition (Blue) --
\newtcolorbox{defbox}[2][]{
  enhanced, breakable,
  colback=defbluebg, colframe=defblue, coltitle=white,
  fonttitle=\bfseries\sffamily,
  title={Definition: #2},
  attach boxed title to top left={yshift=-2mm, xshift=4mm},
  boxed title style={colback=defblue, rounded corners},
  arc=4pt, boxrule=1.2pt,
  left=6pt, right=6pt, top=8pt, bottom=6pt,
  #1
}

% -- Theorem / Proposition (Green) --
\newtcolorbox{thmbox}[2][]{
  enhanced, breakable,
  colback=thmbg, colframe=thmgreen, coltitle=white,
  fonttitle=\bfseries\sffamily,
  title={Theorem: #2},
  attach boxed title to top left={yshift=-2mm, xshift=4mm},
  boxed title style={colback=thmgreen, rounded corners},
  arc=4pt, boxrule=1.2pt,
  left=6pt, right=6pt, top=8pt, bottom=6pt,
  #1
}

% -- Proposition (Green variant) --
\newtcolorbox{propbox}[2][]{
  enhanced, breakable,
  colback=thmbg, colframe=thmgreen!70!black, coltitle=white,
  fonttitle=\bfseries\sffamily,
  title={Proposition: #2},
  attach boxed title to top left={yshift=-2mm, xshift=4mm},
  boxed title style={colback=thmgreen!70!black, rounded corners},
  arc=4pt, boxrule=1.2pt,
  left=6pt, right=6pt, top=8pt, bottom=6pt,
  #1
}

% -- Example (Orange/Yellow) --
\newtcolorbox{exbox}[2][]{
  enhanced, breakable,
  colback=exbg, colframe=exorange, coltitle=white,
  fonttitle=\bfseries\sffamily,
  title={Example: #2},
  attach boxed title to top left={yshift=-2mm, xshift=4mm},
  boxed title style={colback=exorange, rounded corners},
  arc=4pt, boxrule=1.2pt,
  left=6pt, right=6pt, top=8pt, bottom=6pt,
  #1
}

% -- Key Insight / Warning (Red) --
\newtcolorbox{warnbox}[2][]{
  enhanced, breakable,
  colback=warnbg, colframe=warnred, coltitle=white,
  fonttitle=\bfseries\sffamily,
  title={Key Insight: #2},
  attach boxed title to top left={yshift=-2mm, xshift=4mm},
  boxed title style={colback=warnred, rounded corners},
  arc=4pt, boxrule=1.2pt,
  left=6pt, right=6pt, top=8pt, bottom=6pt,
  #1
}

% -- Note (Gray-Blue) --
\newtcolorbox{notebox}[1][]{
  enhanced, breakable,
  colback=noticebg, colframe=noticegray!60, coltitle=white,
  fonttitle=\bfseries\sffamily,
  title={Note},
  attach boxed title to top left={yshift=-2mm, xshift=4mm},
  boxed title style={colback=noticegray!60, rounded corners},
  arc=4pt, boxrule=0.8pt,
  left=6pt, right=6pt, top=8pt, bottom=6pt,
  #1
}

% -- Key Takeaway --
\newtcolorbox{takebox}[1][]{
  enhanced, breakable,
  colback=takeawaybg, colframe=taketawayyellow, coltitle=white,
  fonttitle=\bfseries\sffamily,
  title={Key Takeaway},
  attach boxed title to top left={yshift=-2mm, xshift=4mm},
  boxed title style={colback=taketawayyellow, rounded corners},
  arc=4pt, boxrule=1.2pt,
  left=6pt, right=6pt, top=8pt, bottom=6pt,
  #1
}

% -- Intuition --
\newtcolorbox{intubox}[1][]{
  enhanced, breakable,
  colback=accentpurple!8, colframe=accentpurple!60, coltitle=white,
  fonttitle=\bfseries\sffamily,
  title={Intuition},
  attach boxed title to top left={yshift=-2mm, xshift=4mm},
  boxed title style={colback=accentpurple!80, rounded corners},
  arc=4pt, boxrule=0.8pt,
  left=6pt, right=6pt, top=8pt, bottom=6pt,
  #1
}

% ============================================================
%  PAYOFF MATRIX HELPER
%  Usage: inside tabular, use \payoff{a}{b} for (a,b)
% ============================================================
\newcommand{\payoff}[2]{$(#1,\;#2)$}
\newcommand{\nash}[2]{\underline{$#1$},\underline{$#2$}}  % underline NE payoffs
\newcommand{\Nset}{\mathbb{N}}
\newcommand{\Rset}{\mathbb{R}}
\newcommand{\Sset}{\mathcal{S}}
\newcommand{\Ui}{u_i}

% ============================================================
%  SECTION FORMATTING
% ============================================================
\titleformat{\section}
  {\large\bfseries\sffamily\color{defblue}}
  {\thesection}{1em}{}[\titlerule]

\titleformat{\subsection}
  {\normalsize\bfseries\sffamily\color{thmgreen!80!black}}
  {\thesubsection}{1em}{}

\titleformat{\subsubsection}
  {\normalsize\itshape\sffamily\color{exorange!80!black}}
  {\thesubsubsection}{1em}{}

% ============================================================
%  HEADER / FOOTER
% ============================================================
\pagestyle{fancy}
\fancyhf{}
\fancyhead[L]{\sffamily\small\color{defblue} MTH636 -- Game Theory}
\fancyhead[R]{\sffamily\small\color{defblue} \leftmark}
\fancyfoot[C]{\sffamily\small\thepage}
\renewcommand{\headrulewidth}{0.4pt}

% ============================================================
%  HYPERREF SETUP
% ============================================================
\hypersetup{
  colorlinks=true,
  linkcolor=defblue,
  urlcolor=thmgreen,
  citecolor=accentpurple,
  bookmarksnumbered=true,
  pdftitle={MTH636 Game Theory Notes},
  pdfauthor={MTH636 Student Notes}
}

% ============================================================
%  MISC HELPERS
% ============================================================
\newcommand{\BR}{\text{BR}}        % Best Response
\newcommand{\NE}{\text{NE}}        % Nash Equilibrium
\newcommand{\SPNE}{\text{SPNE}}    % Subgame Perfect NE
\newcommand{\Dom}{\text{dom}}      % Dominance
\newcommand{\mixs}[1]{\sigma_{#1}} % Mixed strategy
\newcommand{\pures}[1]{s_{#1}}     % Pure strategy


% ============================================================
\begin{document}

% Title
\begin{center}
  {\LARGE\bfseries\sffamily\color{defblue} MTH636 -- Game Theory}\\[6pt]
  {\Large\sffamily Lecture 15: Nash's Theorem Proof \& Extensive Form Games}\\[4pt]
\end{center}
\vspace{0.5cm}
\hrule height 1.5pt
\vspace{0.3cm}

% ============================================================
\section{Proof of Nash's Theorem (Continued)}

\begin{thmbox}{Nash's Existence Theorem}
  For every finite game, there is a Nash Equilibrium in mixed strategies.
\end{thmbox}

\textbf{Proof (Contd.):}
Let $G = \langle N, (S_i)_{i \in N}, (u_i)_{i \in N} \rangle$. Let $|N| = n$ and $|S_i| = m_i$ for all $i = 1, \dots, n$.

Recall the gain function $g_i^j : \Sigma \to \mathbb{R}$ defined for each player $i$ and pure strategy $s_i^j$:
\[
  g_i^j(\sigma) = \max \left\{ 0, U_i(s_i^j, \sigma_{-i}) - U_i(\sigma) \right\}
\]
\textbf{Claim 1:} $g_i^j$ is a continuous function of $\sigma$ for all $i$ and $j$. \\
\textbf{Claim 2:} $\sigma$ is an NE if and only if $g_i^j(\sigma) = 0$ for all $i=1,\dots,n$ and $j=1,\dots,m_i$.

\textbf{Our Goal:} We want to construct a mapping $f: \Sigma \to \Sigma$ such that:
\begin{itemize}
    \item $f$ is continuous.
    \item Every fixed point of $f$ is an NE of the game.
\end{itemize}
Then by Brouwer's Fixed Point Theorem, $f$ will have a fixed point, proving an NE exists.

\textbf{Construction of $f$:}
Define $f = (f_1, \dots, f_n)$ where $f_i = (f_i^1, \dots, f_i^{m_i})$ is given by:
\[
  f_i^j(\sigma) = \frac{\sigma_i(s_i^j) + g_i^j(\sigma)}{1 + \sum_{k=1}^{m_i} g_i^k(\sigma)} \quad \forall i=1,\dots,n, \ \forall j=1,\dots,m_i
\]

\textbf{Claim 3:} $f(\sigma) \in \Sigma$.
\begin{proof}
  We must verify two properties for $f_i(\sigma)$ to be a valid probability distribution:
  \begin{enumerate}
      \item $f_i^j(\sigma) \geq 0$: This holds because $\sigma_i(s_i^j) \geq 0$ and $g_i^j(\sigma) \geq 0$.
      \item $\sum_{j=1}^{m_i} f_i^j(\sigma) = 1$: 
      \[ \sum_{j=1}^{m_i} \frac{\sigma_i(s_i^j) + g_i^j(\sigma)}{1 + \sum_{k=1}^{m_i} g_i^k(\sigma)} = \frac{\sum_{j=1}^{m_i} \sigma_i(s_i^j) + \sum_{j=1}^{m_i} g_i^j(\sigma)}{1 + \sum_{k=1}^{m_i} g_i^k(\sigma)} = \frac{1 + \sum_{j=1}^{m_i} g_i^j(\sigma)}{1 + \sum_{k=1}^{m_i} g_i^k(\sigma)} = 1 \]
  \end{enumerate}
\end{proof}

\textbf{Claim 4:} $f$ is a continuous function of $\sigma$. \\
\textit{Proof:} This follows directly since the components $g_i^j(\sigma)$ are continuous functions.

\textbf{Claim 5:} Let $\sigma$ be a fixed point of $f$. Then $g_i^j(\sigma) = \sigma_i(s_i^j) \sum_{k=1}^{m_i} g_i^k(\sigma)$.
\begin{proof}
  Suppose $\sigma$ is a fixed point, meaning $f_i^j(\sigma) = \sigma_i(s_i^j)$. Then:
  \[ \frac{\sigma_i(s_i^j) + g_i^j(\sigma)}{1 + \sum_{k=1}^{m_i} g_i^k(\sigma)} = \sigma_i(s_i^j) \]
  \[ \sigma_i(s_i^j) + g_i^j(\sigma) = \sigma_i(s_i^j) + \sigma_i(s_i^j) \sum_{k=1}^{m_i} g_i^k(\sigma) \]
  \[ \implies g_i^j(\sigma) = \sigma_i(s_i^j) \sum_{k=1}^{m_i} g_i^k(\sigma) \]
\end{proof}

\textbf{Claim 6:} If $\sigma$ is a fixed point, then $\sigma$ is an NE of $G$.
\begin{proof}
  Assume for contradiction that $\sigma$ is a fixed point but \textbf{not} an NE.
  This implies there exists some player $p \in N$ and strategy $s_p^k \in S_p$ such that $g_p^k(\sigma) > 0$.
  Consequently, the sum of gains for player $p$ is strictly positive: $\sum_{k=1}^{m_p} g_p^k(\sigma) > 0$.
  
  As $\sigma$ is a fixed point, by Claim 5:
  \[ g_p^j(\sigma) = \sigma_p(s_p^j) \sum_{k=1}^{m_p} g_p^k(\sigma) \]
  Since the sum is strictly positive, we have:
  \[ g_p^j(\sigma) > 0 \iff \sigma_p(s_p^j) > 0 \quad \forall j=1,\dots,m_p \]
  
  Now consider Player $p$'s expected utility:
  \[ U_p(\sigma) = \sum_{k=1}^{m_p} \sigma_p(s_p^k) U(s_p^k, \sigma_{-p}) \]
  Subtract $U_p(\sigma)$ from both sides (using the fact that $\sum \sigma_p(s_p^k) = 1$):
  \[ 0 = \sum_{k=1}^{m_p} \sigma_p(s_p^k) \left( U(s_p^k, \sigma_{-p}) - U_p(\sigma) \right) \]
  Since $g_p^k(\sigma) > 0$ for any $s_p^k$ in the support, we know $U(s_p^k, \sigma_{-p}) - U_p(\sigma)$ is exactly equal to $g_p^k(\sigma)$. So:
  \[ \sum_{s_p^k \mid \sigma_p(s_p^k) > 0} \sigma_p(s_p^k) g_p^k(\sigma) = 0 \]
  However, since $\sigma_p(s_p^k) > 0 \implies g_p^k(\sigma) > 0$, we are summing strictly positive terms, which cannot sum to 0. This is a contradiction. 
  
  Thus, $\sigma$ must be a Nash Equilibrium. By Brouwer's theorem, such a fixed point exists, so a mixed NE exists.
\end{proof}

% ============================================================
\section{Extensive Form Games: Introduction}

We now model games where players take turns sequentially.
\begin{exbox}{Sequential Battle of the Sexes}
  The man takes his action first, and then the woman sees his choice and makes her own move.
  
  \textbf{Game Tree:}
  \begin{itemize}
      \item Root node (M): Man chooses $F$ (Football) or $D$ (Dance).
      \item If M chooses $F$, we reach node $W_1$. Woman chooses $F$ or $D$. Payoffs: $(2,1)$ for $FF$, $(0,0)$ for $FD$.
      \item If M chooses $D$, we reach node $W_2$. Woman chooses $F$ or $D$. Payoffs: $(0,0)$ for $DF$, $(1,2)$ for $DD$.
  \end{itemize}
  
  \textbf{Strategies:}
  A strategy must specify an action for \textit{every} node a player controls.
  \begin{itemize}
      \item Man's strategies: $\{F, D\}$
      \item Woman's strategies: Specify action at $W_1$ and action at $W_2$. E.g., $FD$ means "F if he plays F, D if he plays D". 
      Total strategies: $\{FF, FD, DF, DD\}$.
  \end{itemize}
  
  \textbf{Normal Form Representation:}
  \begin{center}
  \renewcommand{\arraystretch}{1.4}
  \begin{tabular}{c|cccc}
    $M \setminus W$ & FF & FD & DF & DD \\
    \hline
    F & \payoff{2}{1} & \payoff{2}{1} & \payoff{0}{0} & \payoff{0}{0} \\
    D & \payoff{0}{0} & \payoff{1}{2} & \payoff{0}{0} & \payoff{1}{2} \\
  \end{tabular}
  \end{center}
\end{exbox}

\end{document}
